\section{Related work}\label{sec:related}

\paragraph{Eviction with importance scores.}
H2O formalises KV eviction as keeping a balance of recent tokens and
\emph{heavy hitters} with the largest cumulative attention mass, and gives
a greedy-retention argument in the language of submodular
functions~\cite{zhang2023h2o}. ScissorHands advances the
persistence-of-importance hypothesis---pivotal tokens tend to stay
pivotal---and manages the cache accordingly~\cite{liu2023scissorhands}.
TOVA drops, at each step, the single token with the lowest current
attention, letting the cache size follow the attention distribution
\cite{oren2024tova}; in our comparison we cap such policies at the
contract for comparability, following the convention of reporting capped
variants explicitly. StreamingLLM shows that keeping a small prefix of
\emph{attention sinks} plus a recency window stabilises streaming
generation indefinitely~\cite{xiao2024streamingllm}; sinks have since
become standard protected structure, and we adopt them in every scoring
policy. Our \acs{} score is an exponentially decayed cumulative mass: it
inherits the accumulation of H2O but discounts stale mass, and
Proposition~\ref{prop:staleness} gives a clean separation where the
decay strictly matters. Decay-weighted importance is, by itself, an
incremental step; our contribution is to make it part of a stack whose
memory behaviour is provably consistent.

\paragraph{Layer-wise budget allocation.}
PyramidKV popularised non-uniform layer budgets with a static pyramidal
profile---more cache for lower layers, where attention is
dispersed---matching full-cache performance on LongBench at $12\%$
retention~\cite{cai2024pyramidkv}. CAKE frames eviction as cake slicing:
layer preferences are assessed from spatial and temporal attention
dynamics and memory is managed in a cascading manner, with an eviction
indicator that accounts for shifting importance over
time~\cite{qin2025cake}. LAVa derives a layer-attention output loss
metric that enables dynamic head- and layer-wise budgets without
training~\cite{shen2025lava}. These methods supply strong empirical
evidence that allocation matters, but their budget rules are heuristic:
the profile shapes are fixed priors, and the dynamic rules are validated
by end-task accuracy rather than by any guarantee on the realised cache
size. Our \laba{} mechanism differs in two ways. First, the allocation is
computed at every step as the projection of measured dispersion weights
onto the exact budget set (Theorem~\ref{thm:exact}), so the contract holds
by construction rather than by tuning. Second, because the projection
accepts arbitrary weight vectors, the same machinery instantiates uniform
budgets (H2O-style), static pyramids (PyramidKV-style), and our adaptive
weights, which is what makes the matched-entry comparison of
Section~\ref{sec:exp} possible. Very recent systems work continues to
expand this direction, which we view as further evidence that allocation
is a first-class design axis.

\paragraph{Query-aware and prefill-time selection.}
A complementary family avoids the history-only weakness of eviction
scores. SnapKV compresses at prefill time by observing which positions a
lookahead window of prompts attends to~\cite{li2024snapkv}. Quest keeps
the full cache resident but loads only pages whose key ranges can
interact strongly with the current query, using page-level min/max
statistics~\cite{tang2024quest}. DuoAttention distinguishes retrieval
heads from streaming heads with a learned mask and keeps full caches only
for the former~\cite{xiao2024duoattention}. These methods motivate the
needle-retention diagnostics we report in Section~\ref{sec:exp}: when a
history-only policy evicts a needle, no amount of decoding-time cleverness
recovers it, and measuring retention in the cache isolates that failure
mode from generation errors.

\paragraph{Quantisation and systems.}
KIVI quantises keys per-channel and values per-token to 2 bits with
tuning-free rounding, attacking the same memory wall
orthogonally~\cite{liu2024kivi}. PagedAttention manages cache memory at
page granularity~\cite{kwon2023paged}, and recent systems work notes that
kernel-level constraints---FlashAttention never materialising attention
scores, paged allocators reclaiming memory only at page
granularity---shape what eviction policies can actually be
deployed~\cite{dao2023flashattention2,mao2026infra}. Two of our
implementation decisions respond directly to these constraints: we extract
scores from eager attention kernels (so the accounting of scoring cost is
explicit and honest), and our contract is defined over \emph{entry
counts}, which upper-bounds what any page-based allocator would need to
hold. Robustness studies of compression, including security and
instruction-following side effects, further motivate reporting
multi-instruction behaviour rather than perplexity
alone~\cite{ma2026safety}.

\paragraph{Position versus the state of the art.}
Compared with the strongest layer-adaptive methods (CAKE, LAVa), \sage{}
does not claim better end-task accuracy at scale; our claims are
orthogonal and, we argue, prior in importance: a provably exact global
budget with protected floors, an analysis that explains \emph{when}
decayed scoring and dispersion weighting help, and a matched-entry
protocol that makes cross-method comparisons meaningful. The comparison in
Section~\ref{sec:exp} is therefore deliberately internal---same harness,
same projector, same inputs, same accounting---and its scope (models to
$0.5$B, contexts to $4$K, CPU) is reported as a first-class limitation
rather than hidden.
